\chapter{Proposed System Architecture}
\label{chap:chapter3}

\section{System Overview}
\label{sec:system_overview}

The proposed system is designed to provide automatic scoring of ISSF paper targets
using a low-cost USB camera and computer vision techniques. In addition to shot detection
and scoring, the system supports data storage and visualization of a virtual target to enable the management of
athletes, training sessions and competitions.

At a high level, the system consists of five main components: image acquisition, image
processing, scoring logic, data management and visualization. A fixed USB camera captures
images of the paper target after each shot and transmits them to a personal computer for
processing. The image-processing subsystem analyzes the acquired images to localize the target,
detect bullet holes and estimate their geometric properties. Based on these measurements, the
scoring module computes ISSF-compliant decimal scores. The resulting coordinates and scores are
stored in a database, where they are associated with athletes, sessions and competitions, and
can be visualized through a virtual target representation.

Figure~\ref{fig:system_architecture} illustrates the high-level architecture of the proposed
system and the main processing stages from image acquisition to target visualization.

\begin{figure}[htbp]
  \centering
  \includegraphics[
    width=0.9\linewidth,
    height=0.45\textheight,
    keepaspectratio
  ]{system_architecture.png}
  \caption{High-level architecture of the proposed system.}
  \label{fig:system_architecture}
\end{figure}


While the complete system provides end-to-end functionality for automatic target scoring, the
primary focus of this thesis is the image-processing subsystem. In particular, the research
addresses the problem of robust and accurate bullet-hole localization under the constraints
identified in Chapter~2. The remaining components, including data storage and visualization,
serve to integrate the image-processing results into a usable scoring system but are not the
main object of investigation.

The image-processing subsystem is designed as a modular pipeline, allowing different detection
and refinement strategies to be evaluated and compared. Robust bullet-hole detection is treated
separately from precise geometric refinement, enabling the system to combine complementary
approaches and to merge robustness and accuracy. The following sections describe the design
of this subsystem in detail, beginning with image acquisition assumptions and progressing
through detection, refinement and scoring.


\section{Image Acquisition and Operating Assumptions}
\label{sec:image_acquisition}

Images of the paper target are acquired using a fixed USB camera connected to a personal
computer. The camera is mounted at a known and approximately constant distance from the
target and remains stationary during operation. This fixed setup allows the use of geometric
constraints and calibration procedures in later processing stages. 

Figure~\ref{fig:camera_scheme} shows that the camera will be placed at an angle, looking
at the targets from the side of the shooting line.

\begin{figure}[htbp]
  \centering
  \includegraphics[
    width=0.85\linewidth,
    height=0.45\textheight,
    keepaspectratio
  ]{camera_scheme.png}
  \caption{Camera placement relative to the shooting line and paper target.}
  \label{fig:camera_scheme}
\end{figure}


The system is intended to operate in an indoor shooting range, where lighting conditions can be
controlled to a certain extent but are not assumed to be perfectly uniform. Variations in
illumination, shadows and reflections may still occur. No additional white background board is
assumed behind the target by default.

However, the operating environment allows for limited control of background appearance if
required. In particular, the use of a uniform background color behind the target can be
introduced to make segmentation easier and more robust. This option provides additional flexibility while
remaining consistent with the low-cost and controlled indoor deployment scenario considered in
this work.

Images are captured at a sufficient resolution to allow reliable detection of bullet holes and
target features. Image acquisition is triggered after each shot, either manually or through an
external signal, such that each image contains at most one new bullet hole relative to the
previous state of the target.

These assumptions reflect a realistic low-cost deployment scenario and are consistent with the
constraints discussed in Chapter~2. They guide the design of the image-processing pipeline
described in the following sections.

\section{Image-Processing Subsystem and Considered Methods}
\label{sec:ip_methods}

The image-processing subsystem is structured as a sequence of modular processing stages, as
illustrated in Fig.~\ref{fig:system_architecture}. For each stage, different algorithmic strategies
are considered based on the trade-off between robustness, accuracy and computational
complexity. This section outlines the main methods.

\subsection*{Target Localization}

Accurate localization of the target center is essential, as all subsequent measurements are
defined relative to this reference. Two complementary strategies are considered. A learning-based
approach may be used to detect the target or the central black region using a CNN
trained on target images captured under different conditions, including targets
with existing bullet holes. Alternatively, classical image-processing methods based on intensity
thresholding and geometric constraints may be applied to segment the black region of the target
and estimate its center. In both cases, the initial detection defines a region of interest that can
be further refined using classical methods operating on a cropped and normalized image.

\subsection*{Bullet-Hole Detection (Localization)}

Bullet-hole detection localizes all bullet holes within the target region.
Learning-based object detection methods, including both
pretrained detectors (e.g., YOLO-style or region-based CNNs) and custom lightweight CNNs,
are considered for this stage. To reduce data collection effort and improve
generalization, synthetic data generation may be explored to pretrain networks on circular
shapes resembling bullet holes under varying contrast and noise conditions.

As an alternative or complement to learning-based detection, classical approaches based on
intensity thresholding and morphological operations may be employed, particularly in scenarios
where background appearance can be controlled. Under controlled conditions, these methods are very robust.

\subsection*{Overlap Handling and Temporal Information}

In cases where multiple bullet holes are present in close proximity, a single detected region may
contain more than one impact. To address this situation, segmentation-based techniques such as
watershed algorithms driven by distance transforms are considered for separating overlapping
bullet holes. The number of impacts may be estimated to guide the
application of such methods.

When images are acquired sequentially during live shooting, temporal information can also be
used. Image subtraction between consecutive frames may be used to isolate newly created
bullet holes. In this approach, the camera movement must be approximately zero and shooting timing
must be detected.

\subsection*{Geometric Refinement and Center Estimation}

Following coarse detection, a geometric refinement stage is applied to estimate precise bullet-hole
properties, including center coordinates and hole radius. Classical contour-based methods are
considered for this purpose, relying on thresholding, edge detection and contour extraction within
a cropped region of interest. Circle fitting techniques can then be applied to the extracted
contours to estimate the bullet-hole center.

As an alternative, Hough transform-based circle detection may be used to directly estimate
circle parameters from edge maps. Learning-based refinement approaches, such as small regression
or segmentation networks operating on localized regions, are also considered as potential
alternatives to purely classical methods.

\subsection*{Scoring Integration}

The outputs of the target localization and geometric refinement stages are combined to compute
radial distances between bullet-hole centers and the target center. These distances are then
converted into decimal scores according to the ISSF scoring rules described in Chapter~2. The
resulting shot coordinates and scores are forwarded to the scoring, data management and
visualization components of the system, and a virtual target representation is generated.


